\section*{第 \ref{ch:b3:virology} 章——病毒学}

\begin{solution}{exo:b3:virology:1}
脊髓灰质炎\omterm{def:b3:virology:virus}{病毒}：IV 型，正链 RNA——依赖 RNA 的 RNA 聚合酶。
流感\omterm{def:b3:virology:virus}{病毒}：V 型，负链 RNA——自带 RNA 聚合酶，装于
\omterm{def:b3:virology:virus}{病毒粒子}内。HIV：VI 型——逆转录酶（及整合酶）。单纯疱疹
\omterm{def:b3:virology:virus}{病毒}：I 型，双链 DNA——宿主原则上能够提供
其所需的一切，但它自带 DNA 聚合酶和胸苷激酶。
乙型肝炎\omterm{def:b3:virology:virus}{病毒}：VII 型——逆转录酶，用于把前基因组 RNA
拷贝为 DNA。轮状\omterm{def:b3:virology:virus}{病毒}：III 型，双链 RNA——依赖 RNA 的 RNA
聚合酶位于\omterm{def:b3:virology:virus}{病毒粒子}内。
\end{solution}

\begin{solution}{exo:b3:virology:2}
噬菌体的 DNA（$^{32}$P）进入细菌，蛋白质
（$^{35}$S）留在外面，可被剪切脱落；子代噬菌体
继承了 $^{32}$P。由此可知，DNA 是噬菌体注入的物质，也是
指导新噬菌体合成的物质。若蛋白质携带遗传指令，
则 $^{35}$S 会在细菌内被检出，并传给子代。
\end{solution}

\begin{solution}{exo:b3:virology:3}
$84/(0.1\times 10^{-7}) = 8.4\times 10^{9}$ 个噬菌斑形成单位每毫升。
显微镜计数每一个颗粒，包括空\omterm{def:b3:virology:virus}{衣壳}、带有缺陷基因组的颗粒，
以及无法附着或注入的颗粒；只有
完成感染的颗粒才能形成噬菌斑。
\end{solution}

\begin{solution}{exo:b3:virology:4}
附着：进入抑制剂（maraviroc），中和抗体。
进入与融合：融合抑制剂，\omterm{def:b3:membrane-traffic:endocytosis}{内体}酸化阻断剂。
脱壳：结合\omterm{def:b3:virology:virus}{衣壳}的药物（pleconaril；lenacapavir）。表达
与复制：\omterm{def:b3:virology:drugs}{核苷类似物}，逆转录酶和
整合酶抑制剂，\omterm{prop:b3:virology:pathogenesis}{干扰素}，RNA 干扰。\omterm{def:b3:genomics:assembly}{组装}：\omterm{def:b3:virology:drugs}{蛋白酶
抑制剂}（未被切割的多聚蛋白无法\omterm{def:b3:genomics:assembly}{组装}）。释放：
神经氨酸酶抑制剂（流感），以及细胞毒性 T 细胞
在\omterm{def:b3:virology:virus}{病毒}释放前杀死宿主细胞。
\end{solution}

\begin{solution}{exo:b3:virology:5}
一个 \qty{30}{kDa} 的蛋白质含 $30\,000/110 = 273$ 个残基，对应 $818$
个核苷酸：\qty{0.82}{kb}，占 \qty{4.5}{kb} 基因组的 \qty{18}{\%}。
若分别编码 $180$ 个拷贝，需 \qty{147}{kb}，
为整个基因组的三十三倍；\omterm{def:b3:virology:virus}{衣壳}蛋白本身
（\qty{5.4}{MDa}）的质量是基因组（\qty{1.5}{MDa}）的四倍。
\end{solution}

\begin{solution}{exo:b3:virology:6}
$V = 10^{5}(3e^{-0.5t} - 0.5e^{-3t})/2.5$：$t = 0.5$ 时 $8.9\times
10^{4}$；$t = 2$ 时 $4.4\times 10^{4}$；$t = 10$ 时 $8\times 10^{2}$。低于
$50$ 的时刻满足 $1.2\times 10^{5}e^{-0.5t} = 50$，即 $t = 2\ln(2400) \approx
\qty{16}{d}$。
\end{solution}

\begin{solution}{exo:b3:virology:7}
$uL = 0.36$；无误复制的比例为 $e^{-0.36} = 0.70$。三倍长度时：$uL =
1.08$，$e^{-1.08} = 0.34$——多数子代现在携带突变，最适
序列被稀释至阈值附近。一种基因组长为
\qty{30}{kb}，$u = 10^{-4}$ 时，$uL = 3$，仅 \qty{5}{\%}
的拷贝无误，已超过阈值；其外切酶把 $u$ 降至
$10^{-5}$，$uL = 0.3$，无误拷贝占四分之三。
\end{solution}

\begin{solution}{exo:b3:virology:8}
漂移：血凝素和神经氨酸酶中的点突变
在抗体选择下累积，因此每个季节的毒株
部分是新的，去年的免疫力也部分失效。转变：人类毒株与动物毒株
共同感染同一细胞，使八个节段
发生\omterm{def:b3:virology:evolution}{重配}，来自禽类或猪的血凝素亚型出现在
适应人类的\omterm{def:b3:virology:virus}{病毒}中。人群中无人对新亚型具有抗体，因此
整个人群同时易感——这正是
大流行的条件，1918 年、1957 年、1968 年和 2009 年
每次都由一种新的血凝素引发。
\end{solution}

\begin{solution}{exo:b3:virology:9}
疱疹\omterm{def:b3:virology:virus}{病毒}胸苷激酶使阿昔洛韦磷酸化，这一过程
细胞激酶几乎不作用；随后细胞激酶把它
完成为三磷酸形式，\omterm{def:b3:virology:virus}{病毒} DNA 聚合酶将其掺入，由于缺少
3$'$ 羟基，使链延伸终止。未感染的细胞从不产生
活性形式，\omterm{def:b3:virology:virus}{病毒}聚合酶比细胞聚合酶更偏好它。
耐药性：\omterm{def:b3:virology:virus}{病毒}胸苷激酶丢失或改变（最常见），
或\omterm{def:b3:virology:virus}{病毒}聚合酶发生突变，使其排斥
该类似物。
\end{solution}

\begin{solution}{exo:b3:virology:10}
当 $T$ 恒定时，方程关于 $(I, V)$ 为线性，矩阵为
$\begin{pmatrix} -\delta & \beta T \\ p & -c \end{pmatrix}$，迹
$-(\delta + c) < 0$，行列式为 $\delta c - \beta Tp$。
本征值的迹为负，因此仅当行列式为负时，其中之一为正：
即 $\beta Tp > c\delta$，亦即 $R_{0} = \beta
Tp/(c\delta) > 1$。因子含义：$p/\delta$ 为一个受感染细胞
一生中产生的\omterm{def:b3:virology:virus}{病毒粒子}数；$\beta T/c$ 为一个\omterm{def:b3:virology:virus}{病毒粒子}
一生中感染的细胞数；二者之积即一个受感染细胞
所产生的受感染细胞数。
\end{solution}

\begin{solution}{exo:b3:virology:11}
$10^{6} = 2^{20}$：二十个半衰期，即 $880$ 个月，约 $73$ 年——
超过一个人的寿命。阻断复制的药物不影响
未被转录的前\omterm{def:b3:virology:virus}{病毒}；储库持续存在，停药时
只需一个重新激活的细胞即可重新引发感染。彻底治愈
必须清除或永久沉默储库：杀死潜伏细胞
（在治疗下使其重新激活，从而死亡或被杀死），切除或
使前\omterm{def:b3:virology:virus}{病毒}失活，或用\omterm{def:b3:virology:virus}{病毒}无法感染的免疫系统
取代原有免疫系统。
\end{solution}

\begin{solution}{exo:b3:virology:12}
$uL = 0.03$ 个突变每基因组；$10^{9}$ 个基因组每天携带 $3\times
10^{7}$ 个突变，分布于 $9\times 10^{4}$ 种可能的单突变中：每种
突变每天约出现 $300$ 次。某一特定双突变体的出现率为
每基因组 $10^{-12}$，即每天 $10^{-3}$——约三年一次；三突变体
的出现率为 $10^{-18}$，实际上从不出现。两种药物若耐药突变彼此独立，
原则上即已足够，三种药物则可为依从性差留出余地。
两种药物结合同一种酶时，可能被一个突变同时击败，
该突变改变了位点或酶的构象，二者的耐药突变
并不独立，因此只能计为少于两种。
\end{solution}

\begin{solution}{pb:b3:virology:1}
\textbf{1.} $2\times 10^{5}\times 1.5\times 10^{4} = 3\times 10^{9}$
\omterm{def:b3:virology:virus}{病毒粒子}。
\textbf{2.} 每天清除 $cV = 3\times 3\times 10^{9} \approx 10^{10}$ 个；
产生量同样如此。
\textbf{3.} $I^{*} = cV^{*}/p = 9\times 10^{9}/500 = 1.8\times 10^{7}$
感染细胞；每天死亡 $\delta I^{*} = 9\times 10^{6}$ 个。
\textbf{4.} RNA $2\times 9700\times 330 = 6.4\times 10^{6}$ Da；蛋白质
$2000\times 25\,000 = 5\times 10^{7}$ Da；合计 $5.6\times 10^{7}$ Da
，$\approx 10^{-16}$ g。全部\omterm{def:b3:virology:virus}{病毒粒子}合计：$3\times 10^{9}\times
10^{-16} = 3\times 10^{-7}$ g，约三分之一微克——
比一粒沙子轻三千倍。
\textbf{5.} $1.8\times 10^{7}/10^{12} = 2\times 10^{-5}$，占细胞库的份额。
连续十年每天 $9\times 10^{6}$ 个细胞死亡，合计 $3\times 10^{10}$ 个细胞，
占细胞库的百分之三：细胞库失效并非算术上的
减法，而是因为十年的强制再生与慢性
激活耗尽了补充它们的能力。
\textbf{6.} $9\times 10^{6}/24 \approx 4\times 10^{5}$ 个凋亡细胞残体
任一时刻都在被清除。
\textbf{7.} $V = 2\times 10^{5}(3e^{-0.5t} - 0.5e^{-3t})/2.5$: $t = 1$:
$1.4\times 10^{5}$; $t = 3$: $5.4\times 10^{4}$; $t = 7$: $7\times
10^{3}$; $t = 14$: $2\times 10^{2}$.
\textbf{8.} $2.4\times 10^{5}e^{-0.5t} = 50$: $t = 2\ln(4800) \approx 17$
天。
\textbf{9.} 第 3--10 天位于慢时期：该段 $\ln V$ 的斜率为
$-\delta$（自 $5.4\times 10^{4}$ 降至 $1.6\times 10^{3}$，历时七天，
$\delta = 0.50$）。快时期仅持续数小时（$1/c$），因此确定 $c$ 需要
在第一天每隔数小时取样，拟合到双指数
公式。
\textbf{10.} 潜伏感染的细胞，它们在
重新激活时释放\omterm{def:b3:virology:virus}{病毒}；其衰减速率为 $\ln 2/44$ 每月，约合
$5\times 10^{-4}$ 每天。
\textbf{11.} \omterm{def:b3:virology:virus}{病毒粒子} $\ln 2/3 = \qty{0.23}{d}$，约 \qty{5.5}{h}；感染
细胞 $\ln 2/0.5 = \qty{1.4}{d}$。
\textbf{12.} 恒定的\omterm{thm:b3:virology:dynamics}{病毒载量}看似\omterm{def:b3:virology:virus}{病毒}无所作为，
方程表明它是一种稳态，其中 $10^{10}$ 个\omterm{def:b3:virology:virus}{病毒粒子}被产生
并被清除，每天有 $10^{7}$ 个 T 细胞被感染并杀死。
\textbf{13.} 每基因组 $9700\times 10^{-5} \approx 0.1$ 个突变；
每天 $10^{10}\times 0.1 = 10^{9}$ 个突变。
\textbf{14.} 可能的单突变共 $3\times 9700 \approx \num{29000}$ 种；
每种突变每天出现 $10^{9}/\num{29000} \approx 3\times 10^{4}$ 次。
\textbf{15.} 特定双突变：每基因组 $10^{-10}$，
未治疗时约每天一次；特定三突变：$10^{-15}$，每天 $10^{-5}$ 次——约
三个世纪一次。
\textbf{16.} 每天 $10^{7}\times 10^{-15} = 10^{-8}$ 次，
即每年 $4\times 10^{-6}$ 次。
\textbf{17.} 一种药物停用后，对其余两种药物耐药的突变体
即已足够：特定双突变率为 $10^{-10}$；一周内 $7\times 10^{7}$ 个基因组
给出的期望数为 $7\times 10^{-3}$——仍不太可能，除非
停药期间\omterm{thm:b3:virology:dynamics}{病毒载量}向每天 $10^{10}$ 个回升，届时
约每天一次，逃逸便很可能发生。漏服正是
耐药性出现的途径。
\textbf{18.} 这些类似物共用逆转录酶的活性位点，
而该处一个突变（或一组胸苷类似物相关突变）会改变
酶对其中数种的处理方式，即交叉耐药。两种
核苷类药物只计为少于两种独立药物，因此治疗方案
应组合不同类别：一种核苷类、一种整合酶抑制剂、一种\omterm{def:b3:virology:drugs}{蛋白酶
抑制剂}或一种非核苷类药物。
\textbf{19.} $20$ 个半衰期，$880$ 个月，$73$ 年；降至 $10^{4}$ 时，
$\log_{2}100 = 6.6$ 个半衰期，$290$ 个月，$24$ 年。
\textbf{20.} 由一个\omterm{def:b3:virology:virus}{病毒粒子}增至 $3\times 10^{9}$ 个，设每代 $R_{0} = 8$，
则需 $\ln(3\times 10^{9})/\ln 8 = 10.5$ 代，约三周。
\textbf{21.} 冲击并杀灭：在治疗下使潜伏细胞重新激活，以便
死亡或被杀灭——障碍：没有药物能激活它们
全部，且被激活的细胞也不能被可靠地杀灭。基因编辑：
切除前\omterm{def:b3:virology:virus}{病毒}，或在患者细胞中破坏受体 CCR5；
障碍在于须触达每一个细胞。（以 CCR5 缺陷型供体细胞替换骨髓，
已治愈少数患者，其风险
只在本就需要移植时方可接受。）
\textbf{22.} $1 - 1/R_{0}$：$R_{0} = 3$ 时为 \qty{67}{\%}，
$5$ 时为 \qty{80}{\%}。
\textbf{23.} 治疗使患者体内 $\beta \to 0$，\omterm{thm:b3:virology:dynamics}{病毒载量}
依定理降至几近于零，而传播
与载量成正比，因而停止；对每位感染者进行治疗，使
人群的 $R_{0}$ 降至 1 以下。
\textbf{24.} $f = (1 - 1/R_{0})/E$：取 $R_{0} = 4$、$E = 0.7$，$f =
0.75/0.7 = 1.07$——超过了全部人口：仅靠\omterm{def:b3:virology:drugs}{疫苗}无法达到
阈值；若 $E = 0.9$，则 $f = 0.83$。
\textbf{25.} 每天产生约 $10^{10}$ 个\omterm{def:b3:virology:virus}{病毒粒子}，
约第 $17$ 天即检测不到；治疗期间每年约有 $4\times 10^{-6}$ 个三突变体。
\end{solution}
